% =============================================================================
% Ethics Statement and Broader Impact -- paper-05 fork (2026-10-19 ARS revision)
% Forked from ethics_statement.tex, which ethics_wrapper.tex and the archived U01
% roots still use; edit this file for platform_hybrid/paper/paper.tex only.
% Intended placement: after the Conclusion section, before References.
% Label: \label{sec:impact_standalone}
% Target length: >= 2 pages (NeurIPS Broader Impact guidance).
% =============================================================================

\section{Ethics, Limitations, and Broader Impact}
\label{sec:impact_standalone}

This statement complements Section~\ref{sec:limitations} with a dual-use
analysis, the status of compute and carbon accounting, data provenance, the
closed-source training backend, and methodological limits we did not close.

% ---------------------------------------------------------------------------
\subsection{Dual-Use Analysis: Misuse Risks of Reasoning RL}
\label{sec:dualuse}
% ---------------------------------------------------------------------------

\paragraph{Threat model.}
\ours{} releases (i) training scripts that apply GRPO to the GSM8K
\cite{cobbe2021training} mathematical reasoning benchmark and to the
Salesforce xLAM-Function-Calling-60k \cite{liu2024xlam} tool-use corpus,
(ii) twelve public Hugging Face repositories
\texttt{arvindcr4/tinker-rl-bench-<tag>}, enumerated with their revisions in
\path{platform_hybrid/paper/manifests/p05_hub_repositories.json}; ten hold only
result JSON, step logs, and model cards, and two
(\texttt{ppo\_gsm8k\_Qwen3-8B\_s42}, \texttt{ppo\_gsm8k\_Llama-3.1-8B-Instruct\_s42})
hold LoRA adapter weights from the Modal PPO runs, and (iii) diagnostic code for
Zero-Variance Fraction, reward-stability, and length-bias analysis. We
analyse four concrete misuse pathways these artefacts could, in principle,
enable, and we describe the existing guardrails and the residual risk we
judge acceptable for publication.

\paragraph{M1. Weaponisation of mathematical reasoning.}
The most capability-relevant artefact we release is GRPO training code evaluated
on grade-school arithmetic and multi-step reasoning; the reviewed evidence is
inconclusive about a capability improvement. A malicious
operator could attempt to extend this pipeline to tasks of greater
dual-use concern, e.g., chemistry olympiad problems, cryptographic key
recovery, or financial fraud. We judge the \emph{marginal} uplift provided
by our release to be small for three reasons. First, no capability change is
established: the base Qwen3-8B checkpoint already scores 164/200 on the
held-out GSM8K slice, and the separate five-adapter check
(Section~\ref{sec:heldout}) is inconclusive. Its seed-level difference is not
a bound on capability change, which remains unmeasured within wide uncertainty
(single base pass, no baseline variance); ``inconclusive'' does not establish
``small uplift''.
Second, the public literature already contains GRPO implementations
\cite{shao2024deepseekmath, openrlhf2024, verl2024, vonwerra2022trl} that are at
least as capable. Third, any reasoning uplift is bounded by the rewarder:
GSM8K rewards use exact-match against a human-written gold answer, which
does not generalise to the domains listed above without a new reward
signal, which itself requires expertise and labelled data.

\paragraph{M2. Reward hacking and the long-term alignment risk.}
GRPO, like all policy-gradient RL from a proxy reward, is vulnerable to
reward hacking: the model learns to exploit idiosyncrasies of the reward
function rather than solve the intended task \cite{skalse2022specificationgaming,
krakovna2020specification}. Length and long-horizon
behaviour were not logged uniformly (Section~\ref{sec:limitations}), so reward
hacking is neither demonstrated nor excluded.
Users who apply our scripts to under-specified reward functions in
production settings (e.g., a custom ``helpfulness'' classifier) may observe
models that appear to improve on held-out metrics while silently
optimising against unintended proxies. We include the Zero-Variance
Fraction and reward-stability diagnostics precisely as retrospective summaries; they do not identify policy drift or reward hacking.

\paragraph{M3. Circumventing safety training via distillation or tool-use.}
The tool-use track trains LoRA adapters that teach a base model to emit
schema-valid JSON function calls \cite{liu2024xlam}. A motivated adversary
could in principle use the same pipeline to teach a base model to emit
jailbreak prompts as ``tool calls'' or to invoke an adversarial external
tool. We mitigate this by (a) not releasing adapter weights for the tool-use runs
(the corresponding Hub repositories hold result metadata only) and releasing only
adapters, not merged weights, elsewhere,
while noting that an adapter can still alter safety behavior and does not
guarantee that instruct-model safeguards remain effective; (b) stating here that the tool-use runs used a narrow schema (five
synthetic tools) and were not safety-evaluated for open-ended function
calling; and (c) not releasing any function-calling harness that would
actually execute emitted calls.

\paragraph{M4. Compute-efficiency externalities.}
This audit makes no demonstrated efficiency claim or group-size recommendation. Improved training efficiency is
itself dual-use: it lowers the cost of running adversarial fine-tuning at
scale. We judge this externality to be modest relative to the already-low
cost of fine-tuning a 0.6B--8B model with commodity LoRA recipes
\cite{hu2022lora}, but we flag it explicitly.

\paragraph{Residual risk.}
After weighing M1--M4 we judge that the reproducibility, calibration, and
pedagogical benefits of releasing \ours{} outweigh the residual misuse
risk. We adopt the standard mitigation of releasing adapters (not merged
weights), documenting intended use and out-of-scope use in model cards,
and publishing alongside the diagnostics a researcher would need to detect
reward hacking.

% ---------------------------------------------------------------------------
\subsection{Compute Cost Accounting}
\label{sec:compute_accounting}
% ---------------------------------------------------------------------------

No reconciled spend ledger exists. The historical platform entries (retained,
for audit only, in the shared \path{platform_hybrid/paper/ethics_statement.tex})
are mutually inconsistent: excluding Colab, their numeric ranges sum to USD
147--177 against an earlier stated USD 130--140 total, and cash versus credits,
dates, cohort overlap, and failed-run inclusion are unreconciled. We therefore
report no spend figures and no project total.

\paragraph{Hidden costs.} Beyond direct platform spend, the project
consumed human labour (approximately six student-FTE-months, unpaid) and
infrastructure generously subsidised by PES University's LLM Research
Group (shared access to one A100 node used for TRL baselines and
pre-submission sanity checks). No author received commercial compensation
from Thinking Machines, Modal, or any other vendor mentioned.

\paragraph{Failed-run accounting.} Some spend was consumed by runs that
failed (JWT expiry, API stalls for GPT-OSS-20b and Kimi-K2 before re-run, the
KL-tracking gradient bug); the failed-run subset is not reconciled, so no
amount is stated.

% ---------------------------------------------------------------------------
\subsection{Carbon Footprint Estimation}
\label{sec:carbon}
% ---------------------------------------------------------------------------

No measured energy or emissions data exist for this project. Tinker does not
disclose hardware allocation or GPU-hours, and the historical Modal GPU-hour
figures conflict (18 versus 250 GPU-hours in different summaries). Earlier
versions tabulated conditional scenarios (TDP $\times$ assumed hours $\times$
PUE $\times$ grid intensity, following \citet{patterson2024empirical} and
\citet{strubell2019energy}); because their hours are unreconciled and their
accelerator counts assumed, any summed total would be read as a footprint it is
not, so no carbon figure is reported. A future release needs a single ledger
with dates, cohort IDs, GPU counts, measured versus assumed hours, and failure
status before reporting totals.

\paragraph{Offset and mitigation.}
We did not purchase carbon offsets. Instead, we have (a) published run
result files and step logs (Section~\ref{sec:reproducibility}) so that
downstream users need not re-run the experiments to inspect them; (b) retained local step-level CSV logs so that
learning curves can be inspected without re-training; and (c) listed the failed and quarantined runs, so that
replicators can skip them. We regard reproducibility-by-artefact as a
more durable mitigation than offsets.

% ---------------------------------------------------------------------------
\subsection{Data Provenance}
\label{sec:provenance}
% ---------------------------------------------------------------------------

\ours{} uses only publicly released research datasets. No private data,
personally identifiable information, or licensed proprietary content is
used. No human annotators were employed during this work. We document
each dataset below.

\paragraph{GSM8K \cite{cobbe2021training}.}
Grade-school math word problems (7{,}473 train / 1{,}319 test)
authored by human writers contracted by OpenAI. Released under the
\textbf{MIT License} via \url{https://github.com/openai/grade-school-math}.
We use the standard train/test splits without modification. The canonical
citation is \citet{cobbe2021training}. Known limitations of GSM8K include
gender-neutral but culturally US-centric word problems (names, currencies,
sports) and a moderate rate of human-labelling errors estimated at
\textasciitilde 2\% by \cite{zhang2024gsm1k}. Our held-out evaluation
respects the test/train split.

\paragraph{Salesforce xLAM-Function-Calling-60k \cite{liu2024xlam}.}
60{,}000 function-calling examples generated by Salesforce Research as
training data for the xLAM agentic model family. Released under the
\textbf{CC BY 4.0} license via \url{https://huggingface.co/datasets/Salesforce/xlam-function-calling-60k}.
We use the public release as-is in tool-use tracks of the broader corpus; no
xLAM-based result is displayed in this audit.
Known limitations: xLAM schemas are synthetic and skew toward cleanly-typed
arguments; the dataset under-represents ambiguous, error-handling, and
multi-turn tool-calling. Attribution to Salesforce is required by the
license and is provided in Section~\ref{sec:acknowledgments} and in the
xLAM-derived model cards.

\paragraph{HumanEval \cite{chen2021codex}.}
164 Python programming problems with unit tests, released by OpenAI under
the \textbf{MIT License}. The local restricted sandbox does not implement an unmodified standard
pass@$k$ evaluation; code results are harness-bound reward probes. Known
limitations: small scale, English-language docstrings, single-language
coverage; documented contamination concerns against frontier pretraining
corpora \cite{riddell2024contamination}.

\paragraph{NuminaMath \cite{li2024numinamath}.}
Used by collaborator Mohammad Rafi for the multi-stage
GSM8K+NuminaMath pipeline. Released under \textbf{Apache 2.0}. We use a
downstream checkpoint reported by Rafi; our repository contains only his
scripts and the Apache-licensed derivative weights.

\paragraph{Open-Platypus \cite{lee2023platypus}.}
3{,}000 SFT examples used by collaborator Madhu for Qwen3-8B code
generation warm-up. Open-Platypus is a curated subset released under
\textbf{CC BY-NC 4.0}; the non-commercial restriction is respected in our
release, which is academic-use only.

\paragraph{Synthetic tool-use prompts (authored).}
The reviewed all-wrong cells use ten authored requests over five synthetic tools
(calculator, weather, search, reminder, translate) defined in the candidate
runner source (Section~\ref{sec:evidence-identity}; hash in the reviewed-cell
manifest). An earlier version of this statement described a different five-tool
corpus under \texttt{data/synthetic\_tools/}; no such directory exists in the
repository, and that description is withdrawn. No user data, scraped web
content, or proprietary API traces are present.

\paragraph{No web scraping, no human subjects.}
We did not scrape any website. We did not run any study with human
participants. No IRB review was therefore required. No data that could
reasonably be considered to contain PII, copyrighted content, or sensitive
personal attributes was used at any stage of training, evaluation, or
reward modelling.

% ---------------------------------------------------------------------------
\subsection{Closed-Source Tinker Acknowledgment}
\label{sec:tinker_closed}
% ---------------------------------------------------------------------------

A central tension in \ours{} is that a substantial fraction of our
headline numbers come from a closed-source commercial platform while our
paper simultaneously advocates for reproducibility and platform
independence. We do not resolve this tension by hiding it; we describe
it precisely.

\paragraph{What Tinker is and is not.}
Tinker \cite{thinkingmachines2024tinker} is a managed LLM fine-tuning and
inference service provided by Thinking Machines, Inc. It exposes a Python
SDK (stated elsewhere as \texttt{tinker==0.16.1}; the Hub model cards of the
reviewed cells state v0.18.1, and the executed version is not bound) that accepts custom loss functions
(\texttt{forward\_backward\_custom}), a limited set of optimisers, and
standard LoRA hyperparameters. Local scripts specify reward normalization, baseline subtraction, custom loss,
and requested LoRA/optimizer configuration. Undisclosed server-side details
include hardware placement, scheduling, execution numerics, and some batching
and sampler behavior. These are separate from the inspectable local objective.

\paragraph{What this means for our claims.}
We have not ruled out seed variance, checkpoint/model-size confounding, evaluator
differences, or backend effects. The arithmetic/GSM8K cross-library identity is
unresolved. Tinker results do not establish an upper bound or algorithmic
superiority. verl/OpenRLHF Task~4 dry runs are engineering artifacts, not
measured alternatives. The contribution remains the centered-reward identity,
descriptive audited cells, and reporting contract.

\paragraph{Reproducibility commitments we can make.}
Section~\ref{sec:reproducibility} distinguishes local source, observed
configuration, undisclosed internals, and evidence access. A dagger denotes only
interrupted execution; backend provenance is stated explicitly. Exact checkpoint
revisions and public access must be verified per artifact; no blanket release
or future rerun commitment is made here.

\paragraph{Key rotation and credential hygiene.}
Our Tinker API key (prefix \texttt{tml-...}) is held in a repository
\texttt{.env.example} placeholder only. The real key is stored in the
authors' password manager and rotated whenever a failure mode suggests
possible exfiltration. No Tinker key is committed to git history in
either the primary repository (\texttt{pes-llm-research/tinker-rl-lab}) or
the mirror (\texttt{arvindcr4/tinker-rl-lab}).

% ---------------------------------------------------------------------------
\subsection{Known Methodological Limits}
\label{sec:methodlimits}
% ---------------------------------------------------------------------------

In addition to the infrastructure failures and platform-specific caveats
enumerated in Section~\ref{sec:limitations}, we flag the following
methodological limits.

\paragraph{Short training horizons.}
Most reviewed Tinker runs have short, often 30-step horizons; this is not
a universal limit for all later scripts or campaigns. This is a
deliberate cost-control choice and means our Tinker numbers are
\emph{early-training snapshots}. We cannot rule out that longer training
would change the qualitative picture. The separate held-out GSM8K check (Section~\ref{sec:heldout}) is a
one-sample seed-level test against a fixed baseline, not paired-seed and not
reviewed-record evidence; it establishes no held-out gain.

\paragraph{Single-seed Tinker experiments.}
Each Tinker configuration in the ledger was run once, except Qwen3-8B-Base,
whose two executions at the same visible settings differ by 13.1 points in
last-10 reward; best practice prescribes 3--5 runs \cite{henderson2018deep}. Tinker results therefore lack
variance estimates and do not support significance testing; we report
them descriptively. Seed-level inference requires the actual independent seed observations and
an identified estimand; reconstructed seed vectors are not eligible.

\paragraph{Train-set reward as primary Tinker metric.}
Tinker runs primarily report reward on training prompts. Only the GSM8K
track was followed up with a separate 200-example held-out evaluation
per seed. Other Tinker tracks (tool-use, xLAM) remain train-set-only and
we cannot distinguish memorisation from generalisation for those
settings.

\paragraph{LoRA only; no full fine-tuning.}
The reviewed LLM fine-tuning runs use LoRA \cite{hu2022lora}, with settings
varying by cohort; the Wave~6 rank sweep also includes rank 4. We have
\emph{not} tested whether full fine-tuning would re-order the
library/algorithm comparisons. The LoRA constraint is practical (cost) but
it means our claims about PPO vs.\ GRPO and TRL vs.\ Tinker are
LoRA-conditional.

\paragraph{Benchmark coverage is narrow.}
The broader corpus touches GSM8K, MATH-500 (exploratory), HumanEval
(subset), and synthetic/xLAM tool use; this audit displays only GSM8K and the
synthetic tool-use cells. We do not evaluate on MT-Bench,
ArenaHard, safety benchmarks (HarmBench, ToxicChat), or truthfulness
benchmarks (TruthfulQA). Any claim about ``reasoning'' in the paper is
strictly a claim about the covered benchmarks.

\paragraph{No human preference data.}
We use verifiable rewards only (exact-match for math, unit-test for code,
schema-validity for tool-use). We do not study RLHF with human preference
data; findings should not be extrapolated to reward-model--based alignment
without further work.

\paragraph{Closed-source implementation opacity (reiterated).}
Comparisons between Tinker GRPO and TRL GRPO are confounded by the
closed-source nature of Tinker's implementation. See
Section~\ref{sec:tinker_closed} for detail.

\paragraph{Cross-platform hardware confounding.}
Tinker, Modal, Colab, and the PES A100 node use different accelerators,
different memory hierarchies, and different inference/training stacks
(Tinker proprietary, Modal vLLM, Colab PyTorch-native, TRL+Accelerate).
Observed differences in reward dynamics are not purely algorithmic.

\paragraph{Carbon estimates are approximate.}
As discussed in Section~\ref{sec:carbon}, Tinker GPU-hour and TDP are
inferred, not measured. Grid intensity is region-average, not marginal.
No carbon figure is reported for that reason.

\paragraph{Exploratory, not confirmatory.}
This is an exploratory case study, not a pre-registered confirmatory
experiment. We did not pre-register primary hypotheses. The conflicting historical correction families and step-based significance
decisions are withdrawn (Section~\ref{sec:statistical-policy}); exploratory
coefficients do not establish calibrated inference.

\paragraph{Geographic and demographic limits.}
All authors are based at a single institution in Bengaluru, India. Our
choice of benchmarks, prompt phrasings, and evaluation design reflects
that context. Independent replication by groups in other regions, on
non-English tasks, and with non-Qwen / non-Llama families is an open
question.

% =============================================================================
% End of ethics_statement.tex
